% Partie : set-theory
% Chapitre : ordinals
% Section : vn

\documentclass[../../../include/open-logic-section]{subfiles}

\begin{document}

\olfileid[fr]{sth}{ordinals}{vn}
\olsection[La construction de von Neumann]{La construction des ordinaux par von Neumann}

Le \olref[sth][ordinals][iso]{thm:woalwayscomparable} suggère une idée.
Nous pourrions introduire certains objets, appelés \emph{types
d'ordre}, pour tenir lieu des bons ordres. En notant $\ordtype{A, <}$
le type d'ordre du bon ordre $\tuple{A, <}$, nous souhaiterions obtenir
les deux principes suivants :
\begin{align*}
	\ordtype{A, <} = \ordtype{B, \lessdot} &
	\text{ ssi } \ordeq{\tuple{A, <}}{\tuple{B, \lessdot}}\\
	\ordtype{A, <} < \ordtype{B, \lessdot}&
	\text{ ssi }\ordeq{\tuple{A, <}}{\tuple{B_b, \lessdot_b}}\text{ pour un certain }b \in B
\end{align*}
Nous pourrions, de plus, espérer définir ces types d'ordre \emph{comme
certains ensembles}, tout comme nous pouvons définir les nombres
naturels comme certains ensembles.

La façon la plus courante de procéder, que nous suivrons ici, consiste
à définir ces types d'ordre au moyen de certains ensembles bien
ordonnés \emph{canoniques}. Von Neumann a introduit ces ensembles
canoniques le premier :

\begin{defn}
L'ensemble $A$ est \emph{transitif} si et seulement si
$(\forall x \in A)x \subseteq A$. L'ensemble $A$ est alors un
\emph{ordinal} si et seulement s'il est transitif et bien ordonné par
$\in$.
\end{defn}
\noindent
Dans la suite, nous désignerons les ordinaux par des lettres grecques.
La définition entraîne immédiatement que, si $\alpha$ est un ordinal,
alors $\tuple{\alpha, \in_\alpha}$ est un bon ordre, où
$\in_\alpha = \Setabs{\tuple{x, y} \in \alpha^2}{x \in y}$. Ainsi,
au prix d'un léger abus de notation, nous pouvons dire simplement que
$\alpha$ est \emph{lui-même} un bon ordre.

Voici les premiers ordinaux :
\[
	\emptyset, \{\emptyset\},
	\{\emptyset, \{\emptyset\}\},
	\{\emptyset, \{\emptyset\}, \{\emptyset, \{\emptyset\}\}\}, \ldots
\]
On reconnaît les premiers ordinaux rencontrés dans notre axiome de
l'infini, c'est-à-dire dans la définition de $\omega$ donnée par von
Neumann (voir la \olref[sth][z][infinity-again]{sec}). Ce n'est pas un
hasard. La définition des ordinaux de von Neumann considère les
nombres naturels comme des ordinaux, tout en admettant aussi des
ordinaux transfinis.

Comme toujours, nous pouvons demander : ces ensembles \emph{sont-ils}
les ordinaux ? Ou von Neumann nous a-t-il simplement donné des
ensembles que nous pouvons \emph{traiter} comme les ordinaux ? Les
discussions que cette question suscite ressemblent à celles que nous
avons menées dans la \olref[sfr][rel][ref]{sec},
\olref[sfr][arith][ref]{sec},
\olref[sfr][infinite][dedekindsproof]{sec} et
\olref[sth][z][nat]{sec} ; nous n'insisterons donc pas davantage.
Dans la suite, nous emploierons simplement «~les ordinaux~» pour
parler des «~ordinaux de von Neumann~».

\end{document}
